\section{Introduction}
\label{sec:introduction}

Video collections are growing faster than systems can analyze them exhaustively. In practical settings, a user often wants to issue a natural-language query and immediately retrieve the temporal span in a video that matches it. This problem is commonly studied as natural-language video localization or \task. Benchmarks such as \dataset make the task especially realistic because successful retrieval may require reasoning over both visual content and subtitles rather than either modality alone \cite{lei2020tvr}.

This realism comes with a systems cost. Traditional video analytics pipelines could often be optimized around a fixed set of detectors and predicates, but open-ended language queries require more flexible multimodal models. Recent transformer-based approaches have substantially improved temporal localization quality, yet they also increase online computation, memory traffic, and latency because they must score many temporal candidates while processing high-dimensional visual and textual features. The cost becomes even more visible for video, where temporal length compounds the already expensive spatial processing used in modern vision-language models.

This project studies \system, a retrieval-first design for low-latency \task on \dataset. Instead of asking expensive multimodal localization to score the full corpus uniformly, we first use transcript-based retrieval to route each query to a small set of candidate videos and then apply \xml, a cross-modal localizer, on that retained set. We also evaluate several text-only subtitle-window baselines, including \bm, dense retrieval with sentence-transformer embeddings, and their fusion. These baselines provide useful points of comparison for cheap textual retrieval, but they do not account for the main quality gains in our experiments. The main benefit comes from preserving multimodal localization while reducing the candidate set that \xml must score.

The goal of this work is not to propose a completely new localization architecture. The goal is to understand which components provide the best accuracy-latency trade-off under practical constraints, including precomputed video features and limited tolerance for online inference cost. That framing places the project at the intersection of computer vision and systems. The computer vision side concerns robust temporal grounding from natural-language queries, while the systems side concerns how to avoid spending heavy computation on candidates that are unlikely to contain the answer. The central tension is that a text-only retriever is not strong enough to replace cross-modal localization, but it may still be strong enough to decide where that localization should run.

\paragraph{Contributions.}
\begin{itemize}[leftmargin=*]
  \item We present \system, a modular pipeline for low-latency \task that separates cheap subtitle-based retrieval from expensive cross-modal localization.
  \item We evaluate \bm, dense sentence-transformer retrieval, and their fusion as subtitle-search baselines, showing where text-only retrieval helps and where it falls short.
  \item We adapt \xml inference by using dense transcript retrieval as a video-level router before official VCMR scoring.
  \item We evaluate both quality and runtime, including shared temporal compaction as a negative systems ablation, to treat latency as a first-class metric rather than a secondary implementation detail.
\end{itemize}
